\documentclass[11pt]{article}

    \usepackage[breakable]{tcolorbox}
    \usepackage{parskip} % Stop auto-indenting (to mimic markdown behaviour)
    

    % Basic figure setup, for now with no caption control since it's done
    % automatically by Pandoc (which extracts ![](path) syntax from Markdown).
    \usepackage{graphicx}
    % Keep aspect ratio if custom image width or height is specified
    \setkeys{Gin}{keepaspectratio}
    % Maintain compatibility with old templates. Remove in nbconvert 6.0
    \let\Oldincludegraphics\includegraphics
    % Ensure that by default, figures have no caption (until we provide a
    % proper Figure object with a Caption API and a way to capture that
    % in the conversion process - todo).
    \usepackage{caption}
    \DeclareCaptionFormat{nocaption}{}
    \captionsetup{format=nocaption,aboveskip=0pt,belowskip=0pt}

    \usepackage{float}
    \floatplacement{figure}{H} % forces figures to be placed at the correct location
    \usepackage{xcolor} % Allow colors to be defined
    \usepackage{enumerate} % Needed for markdown enumerations to work
    \usepackage{geometry} % Used to adjust the document margins
    \usepackage{amsmath} % Equations
    \usepackage{amssymb} % Equations
    \usepackage{textcomp} % defines textquotesingle
    % Hack from http://tex.stackexchange.com/a/47451/13684:
    \AtBeginDocument{%
        \def\PYZsq{\textquotesingle}% Upright quotes in Pygmentized code
    }
    \usepackage{upquote} % Upright quotes for verbatim code
    \usepackage{eurosym} % defines \euro

    \usepackage{iftex}
    \ifPDFTeX
        \usepackage[T1]{fontenc}
        \IfFileExists{alphabeta.sty}{
              \usepackage{alphabeta}
          }{
              \usepackage[mathletters]{ucs}
              \usepackage[utf8x]{inputenc}
          }
    \else
        \usepackage{fontspec}
        \usepackage{unicode-math}
    \fi

    \usepackage{fancyvrb} % verbatim replacement that allows latex
    \usepackage{grffile} % extends the file name processing of package graphics
                         % to support a larger range
    \makeatletter % fix for old versions of grffile with XeLaTeX
    \@ifpackagelater{grffile}{2019/11/01}
    {
      % Do nothing on new versions
    }
    {
      \def\Gread@@xetex#1{%
        \IfFileExists{"\Gin@base".bb}%
        {\Gread@eps{\Gin@base.bb}}%
        {\Gread@@xetex@aux#1}%
      }
    }
    \makeatother
    \usepackage[Export]{adjustbox} % Used to constrain images to a maximum size
    \adjustboxset{max size={0.9\linewidth}{0.9\paperheight}}

    % The hyperref package gives us a pdf with properly built
    % internal navigation ('pdf bookmarks' for the table of contents,
    % internal cross-reference links, web links for URLs, etc.)
    \usepackage{hyperref}
    % The default LaTeX title has an obnoxious amount of whitespace. By default,
    % titling removes some of it. It also provides customization options.
    \usepackage{titling}
    \usepackage{longtable} % longtable support required by pandoc >1.10
    \usepackage{booktabs}  % table support for pandoc > 1.12.2
    \usepackage{array}     % table support for pandoc >= 2.11.3
    \usepackage{calc}      % table minipage width calculation for pandoc >= 2.11.1
    \usepackage[inline]{enumitem} % IRkernel/repr support (it uses the enumerate* environment)
    \usepackage[normalem]{ulem} % ulem is needed to support strikethroughs (\sout)
                                % normalem makes italics be italics, not underlines
    \usepackage{soul}      % strikethrough (\st) support for pandoc >= 3.0.0
    \usepackage{mathrsfs}
    

    
    % Colors for the hyperref package
    \definecolor{urlcolor}{rgb}{0,.145,.698}
    \definecolor{linkcolor}{rgb}{.71,0.21,0.01}
    \definecolor{citecolor}{rgb}{.12,.54,.11}

    % ANSI colors
    \definecolor{ansi-black}{HTML}{3E424D}
    \definecolor{ansi-black-intense}{HTML}{282C36}
    \definecolor{ansi-red}{HTML}{E75C58}
    \definecolor{ansi-red-intense}{HTML}{B22B31}
    \definecolor{ansi-green}{HTML}{00A250}
    \definecolor{ansi-green-intense}{HTML}{007427}
    \definecolor{ansi-yellow}{HTML}{DDB62B}
    \definecolor{ansi-yellow-intense}{HTML}{B27D12}
    \definecolor{ansi-blue}{HTML}{208FFB}
    \definecolor{ansi-blue-intense}{HTML}{0065CA}
    \definecolor{ansi-magenta}{HTML}{D160C4}
    \definecolor{ansi-magenta-intense}{HTML}{A03196}
    \definecolor{ansi-cyan}{HTML}{60C6C8}
    \definecolor{ansi-cyan-intense}{HTML}{258F8F}
    \definecolor{ansi-white}{HTML}{C5C1B4}
    \definecolor{ansi-white-intense}{HTML}{A1A6B2}
    \definecolor{ansi-default-inverse-fg}{HTML}{FFFFFF}
    \definecolor{ansi-default-inverse-bg}{HTML}{000000}

    % common color for the border for error outputs.
    \definecolor{outerrorbackground}{HTML}{FFDFDF}

    % commands and environments needed by pandoc snippets
    % extracted from the output of `pandoc -s`
    \providecommand{\tightlist}{%
      \setlength{\itemsep}{0pt}\setlength{\parskip}{0pt}}
    \DefineVerbatimEnvironment{Highlighting}{Verbatim}{commandchars=\\\{\}}
    % Add ',fontsize=\small' for more characters per line
    \newenvironment{Shaded}{}{}
    \newcommand{\KeywordTok}[1]{\textcolor[rgb]{0.00,0.44,0.13}{\textbf{{#1}}}}
    \newcommand{\DataTypeTok}[1]{\textcolor[rgb]{0.56,0.13,0.00}{{#1}}}
    \newcommand{\DecValTok}[1]{\textcolor[rgb]{0.25,0.63,0.44}{{#1}}}
    \newcommand{\BaseNTok}[1]{\textcolor[rgb]{0.25,0.63,0.44}{{#1}}}
    \newcommand{\FloatTok}[1]{\textcolor[rgb]{0.25,0.63,0.44}{{#1}}}
    \newcommand{\CharTok}[1]{\textcolor[rgb]{0.25,0.44,0.63}{{#1}}}
    \newcommand{\StringTok}[1]{\textcolor[rgb]{0.25,0.44,0.63}{{#1}}}
    \newcommand{\CommentTok}[1]{\textcolor[rgb]{0.38,0.63,0.69}{\textit{{#1}}}}
    \newcommand{\OtherTok}[1]{\textcolor[rgb]{0.00,0.44,0.13}{{#1}}}
    \newcommand{\AlertTok}[1]{\textcolor[rgb]{1.00,0.00,0.00}{\textbf{{#1}}}}
    \newcommand{\FunctionTok}[1]{\textcolor[rgb]{0.02,0.16,0.49}{{#1}}}
    \newcommand{\RegionMarkerTok}[1]{{#1}}
    \newcommand{\ErrorTok}[1]{\textcolor[rgb]{1.00,0.00,0.00}{\textbf{{#1}}}}
    \newcommand{\NormalTok}[1]{{#1}}

    % Additional commands for more recent versions of Pandoc
    \newcommand{\ConstantTok}[1]{\textcolor[rgb]{0.53,0.00,0.00}{{#1}}}
    \newcommand{\SpecialCharTok}[1]{\textcolor[rgb]{0.25,0.44,0.63}{{#1}}}
    \newcommand{\VerbatimStringTok}[1]{\textcolor[rgb]{0.25,0.44,0.63}{{#1}}}
    \newcommand{\SpecialStringTok}[1]{\textcolor[rgb]{0.73,0.40,0.53}{{#1}}}
    \newcommand{\ImportTok}[1]{{#1}}
    \newcommand{\DocumentationTok}[1]{\textcolor[rgb]{0.73,0.13,0.13}{\textit{{#1}}}}
    \newcommand{\AnnotationTok}[1]{\textcolor[rgb]{0.38,0.63,0.69}{\textbf{\textit{{#1}}}}}
    \newcommand{\CommentVarTok}[1]{\textcolor[rgb]{0.38,0.63,0.69}{\textbf{\textit{{#1}}}}}
    \newcommand{\VariableTok}[1]{\textcolor[rgb]{0.10,0.09,0.49}{{#1}}}
    \newcommand{\ControlFlowTok}[1]{\textcolor[rgb]{0.00,0.44,0.13}{\textbf{{#1}}}}
    \newcommand{\OperatorTok}[1]{\textcolor[rgb]{0.40,0.40,0.40}{{#1}}}
    \newcommand{\BuiltInTok}[1]{{#1}}
    \newcommand{\ExtensionTok}[1]{{#1}}
    \newcommand{\PreprocessorTok}[1]{\textcolor[rgb]{0.74,0.48,0.00}{{#1}}}
    \newcommand{\AttributeTok}[1]{\textcolor[rgb]{0.49,0.56,0.16}{{#1}}}
    \newcommand{\InformationTok}[1]{\textcolor[rgb]{0.38,0.63,0.69}{\textbf{\textit{{#1}}}}}
    \newcommand{\WarningTok}[1]{\textcolor[rgb]{0.38,0.63,0.69}{\textbf{\textit{{#1}}}}}
    \makeatletter
    \newsavebox\pandoc@box
    \newcommand*\pandocbounded[1]{%
      \sbox\pandoc@box{#1}%
      % scaling factors for width and height
      \Gscale@div\@tempa\textheight{\dimexpr\ht\pandoc@box+\dp\pandoc@box\relax}%
      \Gscale@div\@tempb\linewidth{\wd\pandoc@box}%
      % select the smaller of both
      \ifdim\@tempb\p@<\@tempa\p@
        \let\@tempa\@tempb
      \fi
      % scaling accordingly (\@tempa < 1)
      \ifdim\@tempa\p@<\p@
        \scalebox{\@tempa}{\usebox\pandoc@box}%
      % scaling not needed, use as it is
      \else
        \usebox{\pandoc@box}%
      \fi
    }
    \makeatother

    % Define a nice break command that doesn't care if a line doesn't already
    % exist.
    \def\br{\hspace*{\fill} \\* }
    % Math Jax compatibility definitions
    \def\gt{>}
    \def\lt{<}
    \let\Oldtex\TeX
    \let\Oldlatex\LaTeX
    \renewcommand{\TeX}{\textrm{\Oldtex}}
    \renewcommand{\LaTeX}{\textrm{\Oldlatex}}
    % Document parameters
    % Document title
    \title{Notebook}
    
    
    
    
    
    
    
% Pygments definitions
\makeatletter
\def\PY@reset{\let\PY@it=\relax \let\PY@bf=\relax%
    \let\PY@ul=\relax \let\PY@tc=\relax%
    \let\PY@bc=\relax \let\PY@ff=\relax}
\def\PY@tok#1{\csname PY@tok@#1\endcsname}
\def\PY@toks#1+{\ifx\relax#1\empty\else%
    \PY@tok{#1}\expandafter\PY@toks\fi}
\def\PY@do#1{\PY@bc{\PY@tc{\PY@ul{%
    \PY@it{\PY@bf{\PY@ff{#1}}}}}}}
\def\PY#1#2{\PY@reset\PY@toks#1+\relax+\PY@do{#2}}

\@namedef{PY@tok@w}{\def\PY@tc##1{\textcolor[rgb]{0.73,0.73,0.73}{##1}}}
\@namedef{PY@tok@c}{\let\PY@it=\textit\def\PY@tc##1{\textcolor[rgb]{0.24,0.48,0.48}{##1}}}
\@namedef{PY@tok@cp}{\def\PY@tc##1{\textcolor[rgb]{0.61,0.40,0.00}{##1}}}
\@namedef{PY@tok@k}{\let\PY@bf=\textbf\def\PY@tc##1{\textcolor[rgb]{0.00,0.50,0.00}{##1}}}
\@namedef{PY@tok@kp}{\def\PY@tc##1{\textcolor[rgb]{0.00,0.50,0.00}{##1}}}
\@namedef{PY@tok@kt}{\def\PY@tc##1{\textcolor[rgb]{0.69,0.00,0.25}{##1}}}
\@namedef{PY@tok@o}{\def\PY@tc##1{\textcolor[rgb]{0.40,0.40,0.40}{##1}}}
\@namedef{PY@tok@ow}{\let\PY@bf=\textbf\def\PY@tc##1{\textcolor[rgb]{0.67,0.13,1.00}{##1}}}
\@namedef{PY@tok@nb}{\def\PY@tc##1{\textcolor[rgb]{0.00,0.50,0.00}{##1}}}
\@namedef{PY@tok@nf}{\def\PY@tc##1{\textcolor[rgb]{0.00,0.00,1.00}{##1}}}
\@namedef{PY@tok@nc}{\let\PY@bf=\textbf\def\PY@tc##1{\textcolor[rgb]{0.00,0.00,1.00}{##1}}}
\@namedef{PY@tok@nn}{\let\PY@bf=\textbf\def\PY@tc##1{\textcolor[rgb]{0.00,0.00,1.00}{##1}}}
\@namedef{PY@tok@ne}{\let\PY@bf=\textbf\def\PY@tc##1{\textcolor[rgb]{0.80,0.25,0.22}{##1}}}
\@namedef{PY@tok@nv}{\def\PY@tc##1{\textcolor[rgb]{0.10,0.09,0.49}{##1}}}
\@namedef{PY@tok@no}{\def\PY@tc##1{\textcolor[rgb]{0.53,0.00,0.00}{##1}}}
\@namedef{PY@tok@nl}{\def\PY@tc##1{\textcolor[rgb]{0.46,0.46,0.00}{##1}}}
\@namedef{PY@tok@ni}{\let\PY@bf=\textbf\def\PY@tc##1{\textcolor[rgb]{0.44,0.44,0.44}{##1}}}
\@namedef{PY@tok@na}{\def\PY@tc##1{\textcolor[rgb]{0.41,0.47,0.13}{##1}}}
\@namedef{PY@tok@nt}{\let\PY@bf=\textbf\def\PY@tc##1{\textcolor[rgb]{0.00,0.50,0.00}{##1}}}
\@namedef{PY@tok@nd}{\def\PY@tc##1{\textcolor[rgb]{0.67,0.13,1.00}{##1}}}
\@namedef{PY@tok@s}{\def\PY@tc##1{\textcolor[rgb]{0.73,0.13,0.13}{##1}}}
\@namedef{PY@tok@sd}{\let\PY@it=\textit\def\PY@tc##1{\textcolor[rgb]{0.73,0.13,0.13}{##1}}}
\@namedef{PY@tok@si}{\let\PY@bf=\textbf\def\PY@tc##1{\textcolor[rgb]{0.64,0.35,0.47}{##1}}}
\@namedef{PY@tok@se}{\let\PY@bf=\textbf\def\PY@tc##1{\textcolor[rgb]{0.67,0.36,0.12}{##1}}}
\@namedef{PY@tok@sr}{\def\PY@tc##1{\textcolor[rgb]{0.64,0.35,0.47}{##1}}}
\@namedef{PY@tok@ss}{\def\PY@tc##1{\textcolor[rgb]{0.10,0.09,0.49}{##1}}}
\@namedef{PY@tok@sx}{\def\PY@tc##1{\textcolor[rgb]{0.00,0.50,0.00}{##1}}}
\@namedef{PY@tok@m}{\def\PY@tc##1{\textcolor[rgb]{0.40,0.40,0.40}{##1}}}
\@namedef{PY@tok@gh}{\let\PY@bf=\textbf\def\PY@tc##1{\textcolor[rgb]{0.00,0.00,0.50}{##1}}}
\@namedef{PY@tok@gu}{\let\PY@bf=\textbf\def\PY@tc##1{\textcolor[rgb]{0.50,0.00,0.50}{##1}}}
\@namedef{PY@tok@gd}{\def\PY@tc##1{\textcolor[rgb]{0.63,0.00,0.00}{##1}}}
\@namedef{PY@tok@gi}{\def\PY@tc##1{\textcolor[rgb]{0.00,0.52,0.00}{##1}}}
\@namedef{PY@tok@gr}{\def\PY@tc##1{\textcolor[rgb]{0.89,0.00,0.00}{##1}}}
\@namedef{PY@tok@ge}{\let\PY@it=\textit}
\@namedef{PY@tok@gs}{\let\PY@bf=\textbf}
\@namedef{PY@tok@ges}{\let\PY@bf=\textbf\let\PY@it=\textit}
\@namedef{PY@tok@gp}{\let\PY@bf=\textbf\def\PY@tc##1{\textcolor[rgb]{0.00,0.00,0.50}{##1}}}
\@namedef{PY@tok@go}{\def\PY@tc##1{\textcolor[rgb]{0.44,0.44,0.44}{##1}}}
\@namedef{PY@tok@gt}{\def\PY@tc##1{\textcolor[rgb]{0.00,0.27,0.87}{##1}}}
\@namedef{PY@tok@err}{\def\PY@bc##1{{\setlength{\fboxsep}{\string -\fboxrule}\fcolorbox[rgb]{1.00,0.00,0.00}{1,1,1}{\strut ##1}}}}
\@namedef{PY@tok@kc}{\let\PY@bf=\textbf\def\PY@tc##1{\textcolor[rgb]{0.00,0.50,0.00}{##1}}}
\@namedef{PY@tok@kd}{\let\PY@bf=\textbf\def\PY@tc##1{\textcolor[rgb]{0.00,0.50,0.00}{##1}}}
\@namedef{PY@tok@kn}{\let\PY@bf=\textbf\def\PY@tc##1{\textcolor[rgb]{0.00,0.50,0.00}{##1}}}
\@namedef{PY@tok@kr}{\let\PY@bf=\textbf\def\PY@tc##1{\textcolor[rgb]{0.00,0.50,0.00}{##1}}}
\@namedef{PY@tok@bp}{\def\PY@tc##1{\textcolor[rgb]{0.00,0.50,0.00}{##1}}}
\@namedef{PY@tok@fm}{\def\PY@tc##1{\textcolor[rgb]{0.00,0.00,1.00}{##1}}}
\@namedef{PY@tok@vc}{\def\PY@tc##1{\textcolor[rgb]{0.10,0.09,0.49}{##1}}}
\@namedef{PY@tok@vg}{\def\PY@tc##1{\textcolor[rgb]{0.10,0.09,0.49}{##1}}}
\@namedef{PY@tok@vi}{\def\PY@tc##1{\textcolor[rgb]{0.10,0.09,0.49}{##1}}}
\@namedef{PY@tok@vm}{\def\PY@tc##1{\textcolor[rgb]{0.10,0.09,0.49}{##1}}}
\@namedef{PY@tok@sa}{\def\PY@tc##1{\textcolor[rgb]{0.73,0.13,0.13}{##1}}}
\@namedef{PY@tok@sb}{\def\PY@tc##1{\textcolor[rgb]{0.73,0.13,0.13}{##1}}}
\@namedef{PY@tok@sc}{\def\PY@tc##1{\textcolor[rgb]{0.73,0.13,0.13}{##1}}}
\@namedef{PY@tok@dl}{\def\PY@tc##1{\textcolor[rgb]{0.73,0.13,0.13}{##1}}}
\@namedef{PY@tok@s2}{\def\PY@tc##1{\textcolor[rgb]{0.73,0.13,0.13}{##1}}}
\@namedef{PY@tok@sh}{\def\PY@tc##1{\textcolor[rgb]{0.73,0.13,0.13}{##1}}}
\@namedef{PY@tok@s1}{\def\PY@tc##1{\textcolor[rgb]{0.73,0.13,0.13}{##1}}}
\@namedef{PY@tok@mb}{\def\PY@tc##1{\textcolor[rgb]{0.40,0.40,0.40}{##1}}}
\@namedef{PY@tok@mf}{\def\PY@tc##1{\textcolor[rgb]{0.40,0.40,0.40}{##1}}}
\@namedef{PY@tok@mh}{\def\PY@tc##1{\textcolor[rgb]{0.40,0.40,0.40}{##1}}}
\@namedef{PY@tok@mi}{\def\PY@tc##1{\textcolor[rgb]{0.40,0.40,0.40}{##1}}}
\@namedef{PY@tok@il}{\def\PY@tc##1{\textcolor[rgb]{0.40,0.40,0.40}{##1}}}
\@namedef{PY@tok@mo}{\def\PY@tc##1{\textcolor[rgb]{0.40,0.40,0.40}{##1}}}
\@namedef{PY@tok@ch}{\let\PY@it=\textit\def\PY@tc##1{\textcolor[rgb]{0.24,0.48,0.48}{##1}}}
\@namedef{PY@tok@cm}{\let\PY@it=\textit\def\PY@tc##1{\textcolor[rgb]{0.24,0.48,0.48}{##1}}}
\@namedef{PY@tok@cpf}{\let\PY@it=\textit\def\PY@tc##1{\textcolor[rgb]{0.24,0.48,0.48}{##1}}}
\@namedef{PY@tok@c1}{\let\PY@it=\textit\def\PY@tc##1{\textcolor[rgb]{0.24,0.48,0.48}{##1}}}
\@namedef{PY@tok@cs}{\let\PY@it=\textit\def\PY@tc##1{\textcolor[rgb]{0.24,0.48,0.48}{##1}}}

\def\PYZbs{\char`\\}
\def\PYZus{\char`\_}
\def\PYZob{\char`\{}
\def\PYZcb{\char`\}}
\def\PYZca{\char`\^}
\def\PYZam{\char`\&}
\def\PYZlt{\char`\<}
\def\PYZgt{\char`\>}
\def\PYZsh{\char`\#}
\def\PYZpc{\char`\%}
\def\PYZdl{\char`\$}
\def\PYZhy{\char`\-}
\def\PYZsq{\char`\'}
\def\PYZdq{\char`\"}
\def\PYZti{\char`\~}
% for compatibility with earlier versions
\def\PYZat{@}
\def\PYZlb{[}
\def\PYZrb{]}
\makeatother


    % For linebreaks inside Verbatim environment from package fancyvrb.
    \makeatletter
        \newbox\Wrappedcontinuationbox
        \newbox\Wrappedvisiblespacebox
        \newcommand*\Wrappedvisiblespace {\textcolor{red}{\textvisiblespace}}
        \newcommand*\Wrappedcontinuationsymbol {\textcolor{red}{\llap{\tiny$\m@th\hookrightarrow$}}}
        \newcommand*\Wrappedcontinuationindent {3ex }
        \newcommand*\Wrappedafterbreak {\kern\Wrappedcontinuationindent\copy\Wrappedcontinuationbox}
        % Take advantage of the already applied Pygments mark-up to insert
        % potential linebreaks for TeX processing.
        %        {, <, #, %, $, ' and ": go to next line.
        %        _, }, ^, &, >, - and ~: stay at end of broken line.
        % Use of \textquotesingle for straight quote.
        \newcommand*\Wrappedbreaksatspecials {%
            \def\PYGZus{\discretionary{\char`\_}{\Wrappedafterbreak}{\char`\_}}%
            \def\PYGZob{\discretionary{}{\Wrappedafterbreak\char`\{}{\char`\{}}%
            \def\PYGZcb{\discretionary{\char`\}}{\Wrappedafterbreak}{\char`\}}}%
            \def\PYGZca{\discretionary{\char`\^}{\Wrappedafterbreak}{\char`\^}}%
            \def\PYGZam{\discretionary{\char`\&}{\Wrappedafterbreak}{\char`\&}}%
            \def\PYGZlt{\discretionary{}{\Wrappedafterbreak\char`\<}{\char`\<}}%
            \def\PYGZgt{\discretionary{\char`\>}{\Wrappedafterbreak}{\char`\>}}%
            \def\PYGZsh{\discretionary{}{\Wrappedafterbreak\char`\#}{\char`\#}}%
            \def\PYGZpc{\discretionary{}{\Wrappedafterbreak\char`\%}{\char`\%}}%
            \def\PYGZdl{\discretionary{}{\Wrappedafterbreak\char`\$}{\char`\$}}%
            \def\PYGZhy{\discretionary{\char`\-}{\Wrappedafterbreak}{\char`\-}}%
            \def\PYGZsq{\discretionary{}{\Wrappedafterbreak\textquotesingle}{\textquotesingle}}%
            \def\PYGZdq{\discretionary{}{\Wrappedafterbreak\char`\"}{\char`\"}}%
            \def\PYGZti{\discretionary{\char`\~}{\Wrappedafterbreak}{\char`\~}}%
        }
        % Some characters . , ; ? ! / are not pygmentized.
        % This macro makes them "active" and they will insert potential linebreaks
        \newcommand*\Wrappedbreaksatpunct {%
            \lccode`\~`\.\lowercase{\def~}{\discretionary{\hbox{\char`\.}}{\Wrappedafterbreak}{\hbox{\char`\.}}}%
            \lccode`\~`\,\lowercase{\def~}{\discretionary{\hbox{\char`\,}}{\Wrappedafterbreak}{\hbox{\char`\,}}}%
            \lccode`\~`\;\lowercase{\def~}{\discretionary{\hbox{\char`\;}}{\Wrappedafterbreak}{\hbox{\char`\;}}}%
            \lccode`\~`\:\lowercase{\def~}{\discretionary{\hbox{\char`\:}}{\Wrappedafterbreak}{\hbox{\char`\:}}}%
            \lccode`\~`\?\lowercase{\def~}{\discretionary{\hbox{\char`\?}}{\Wrappedafterbreak}{\hbox{\char`\?}}}%
            \lccode`\~`\!\lowercase{\def~}{\discretionary{\hbox{\char`\!}}{\Wrappedafterbreak}{\hbox{\char`\!}}}%
            \lccode`\~`\/\lowercase{\def~}{\discretionary{\hbox{\char`\/}}{\Wrappedafterbreak}{\hbox{\char`\/}}}%
            \catcode`\.\active
            \catcode`\,\active
            \catcode`\;\active
            \catcode`\:\active
            \catcode`\?\active
            \catcode`\!\active
            \catcode`\/\active
            \lccode`\~`\~
        }
    \makeatother

    \let\OriginalVerbatim=\Verbatim
    \makeatletter
    \renewcommand{\Verbatim}[1][1]{%
        %\parskip\z@skip
        \sbox\Wrappedcontinuationbox {\Wrappedcontinuationsymbol}%
        \sbox\Wrappedvisiblespacebox {\FV@SetupFont\Wrappedvisiblespace}%
        \def\FancyVerbFormatLine ##1{\hsize\linewidth
            \vtop{\raggedright\hyphenpenalty\z@\exhyphenpenalty\z@
                \doublehyphendemerits\z@\finalhyphendemerits\z@
                \strut ##1\strut}%
        }%
        % If the linebreak is at a space, the latter will be displayed as visible
        % space at end of first line, and a continuation symbol starts next line.
        % Stretch/shrink are however usually zero for typewriter font.
        \def\FV@Space {%
            \nobreak\hskip\z@ plus\fontdimen3\font minus\fontdimen4\font
            \discretionary{\copy\Wrappedvisiblespacebox}{\Wrappedafterbreak}
            {\kern\fontdimen2\font}%
        }%

        % Allow breaks at special characters using \PYG... macros.
        \Wrappedbreaksatspecials
        % Breaks at punctuation characters . , ; ? ! and / need catcode=\active
        \OriginalVerbatim[#1,codes*=\Wrappedbreaksatpunct]%
    }
    \makeatother

    % Exact colors from NB
    \definecolor{incolor}{HTML}{303F9F}
    \definecolor{outcolor}{HTML}{D84315}
    \definecolor{cellborder}{HTML}{CFCFCF}
    \definecolor{cellbackground}{HTML}{F7F7F7}

    % prompt
    \makeatletter
    \newcommand{\boxspacing}{\kern\kvtcb@left@rule\kern\kvtcb@boxsep}
    \makeatother
    \newcommand{\prompt}[4]{
        {\ttfamily\llap{{\color{#2}[#3]:\hspace{3pt}#4}}\vspace{-\baselineskip}}
    }
    

    
    % Prevent overflowing lines due to hard-to-break entities
    \sloppy
    % Setup hyperref package
    \hypersetup{
      breaklinks=true,  % so long urls are correctly broken across lines
      colorlinks=true,
      urlcolor=urlcolor,
      linkcolor=linkcolor,
      citecolor=citecolor,
      }
    % Slightly bigger margins than the latex defaults
    
    \geometry{verbose,tmargin=1in,bmargin=1in,lmargin=1in,rmargin=1in}
    
    

\begin{document}
    
    \maketitle
    
    

    
    \section{Exercício: Random Forest e
AdaBoost}\label{exercuxedcio-random-forest-e-adaboost}

Use \texttt{load\_breast\_cancer()} do scikit-learn. Neste conjunto,
\texttt{0\ =\ maligno} e \texttt{1\ =\ benigno}. Adote
\texttt{random\_state=42}. Entregue o notebook com código executado e
respostas fundamentadas nos resultados. Trata-se de dados didáticos; a
comparação não constitui validação clínica.

    \subsection{Questão 1 --- Conheça e prepare os
dados}\label{questuxe3o-1-conheuxe7a-e-prepare-os-dados}

\begin{itemize}
\tightlist
\item
  Carregue o conjunto em um \texttt{DataFrame} com os 30 atributos e a
  variável alvo. Informe dimensões, tipos, valores ausentes, linhas
  duplicadas e quantidade por classe.
\item
  Descreva o tratamento necessário (ou explique por que nenhum é
  necessário). Separe treino (60\%), validação (20\%) e teste (20\%) de
  modo estratificado; mantenha o teste reservado para a questão 5. Por
  que preservar a proporção das classes?
\end{itemize}

    \begin{tcolorbox}[breakable, size=fbox, boxrule=1pt, pad at break*=1mm,colback=cellbackground, colframe=cellborder]
\prompt{In}{incolor}{1}{\boxspacing}
\begin{Verbatim}[commandchars=\\\{\}]
\PY{k+kn}{import}\PY{+w}{ }\PY{n+nn}{pandas}\PY{+w}{ }\PY{k}{as}\PY{+w}{ }\PY{n+nn}{pd}
\PY{k+kn}{from}\PY{+w}{ }\PY{n+nn}{sklearn}\PY{n+nn}{.}\PY{n+nn}{datasets}\PY{+w}{ }\PY{k+kn}{import} \PY{n}{load\PYZus{}breast\PYZus{}cancer}
\PY{k+kn}{from}\PY{+w}{ }\PY{n+nn}{sklearn}\PY{n+nn}{.}\PY{n+nn}{model\PYZus{}selection}\PY{+w}{ }\PY{k+kn}{import} \PY{n}{train\PYZus{}test\PYZus{}split}

\PY{n}{data} \PY{o}{=} \PY{n}{load\PYZus{}breast\PYZus{}cancer}\PY{p}{(}\PY{p}{)}
\PY{n}{X} \PY{o}{=} \PY{n}{pd}\PY{o}{.}\PY{n}{DataFrame}\PY{p}{(}\PY{n}{data}\PY{o}{.}\PY{n}{data}\PY{p}{,} \PY{n}{columns}\PY{o}{=}\PY{n}{data}\PY{o}{.}\PY{n}{feature\PYZus{}names}\PY{p}{)}
\PY{n}{y} \PY{o}{=} \PY{n}{pd}\PY{o}{.}\PY{n}{Series}\PY{p}{(}\PY{n}{data}\PY{o}{.}\PY{n}{target}\PY{p}{,} \PY{n}{name}\PY{o}{=}\PY{l+s+s1}{\PYZsq{}}\PY{l+s+s1}{target}\PY{l+s+s1}{\PYZsq{}}\PY{p}{)}
\PY{n}{df} \PY{o}{=} \PY{n}{pd}\PY{o}{.}\PY{n}{concat}\PY{p}{(}\PY{p}{[}\PY{n}{X}\PY{p}{,} \PY{n}{y}\PY{p}{]}\PY{p}{,} \PY{n}{axis}\PY{o}{=}\PY{l+m+mi}{1}\PY{p}{)}

\PY{n+nb}{print}\PY{p}{(}\PY{l+s+s2}{\PYZdq{}}\PY{l+s+s2}{Dimensoes:}\PY{l+s+s2}{\PYZdq{}}\PY{p}{,} \PY{n}{df}\PY{o}{.}\PY{n}{shape}\PY{p}{)}
\PY{n+nb}{print}\PY{p}{(}\PY{l+s+s2}{\PYZdq{}}\PY{l+s+s2}{Tipos de dados:}\PY{l+s+se}{\PYZbs{}n}\PY{l+s+s2}{\PYZdq{}}\PY{p}{,} \PY{n}{df}\PY{o}{.}\PY{n}{dtypes}\PY{o}{.}\PY{n}{value\PYZus{}counts}\PY{p}{(}\PY{p}{)}\PY{p}{)}
\PY{n+nb}{print}\PY{p}{(}\PY{l+s+s2}{\PYZdq{}}\PY{l+s+s2}{Valores ausentes:}\PY{l+s+s2}{\PYZdq{}}\PY{p}{,} \PY{n}{df}\PY{o}{.}\PY{n}{isnull}\PY{p}{(}\PY{p}{)}\PY{o}{.}\PY{n}{sum}\PY{p}{(}\PY{p}{)}\PY{o}{.}\PY{n}{sum}\PY{p}{(}\PY{p}{)}\PY{p}{)}
\PY{n+nb}{print}\PY{p}{(}\PY{l+s+s2}{\PYZdq{}}\PY{l+s+s2}{Linhas duplicadas:}\PY{l+s+s2}{\PYZdq{}}\PY{p}{,} \PY{n}{df}\PY{o}{.}\PY{n}{duplicated}\PY{p}{(}\PY{p}{)}\PY{o}{.}\PY{n}{sum}\PY{p}{(}\PY{p}{)}\PY{p}{)}
\PY{n+nb}{print}\PY{p}{(}\PY{l+s+s2}{\PYZdq{}}\PY{l+s+s2}{Quantidade por classe:}\PY{l+s+se}{\PYZbs{}n}\PY{l+s+s2}{\PYZdq{}}\PY{p}{,} \PY{n}{df}\PY{p}{[}\PY{l+s+s1}{\PYZsq{}}\PY{l+s+s1}{target}\PY{l+s+s1}{\PYZsq{}}\PY{p}{]}\PY{o}{.}\PY{n}{value\PYZus{}counts}\PY{p}{(}\PY{p}{)}\PY{p}{)}

\PY{n}{X\PYZus{}train}\PY{p}{,} \PY{n}{X\PYZus{}temp}\PY{p}{,} \PY{n}{y\PYZus{}train}\PY{p}{,} \PY{n}{y\PYZus{}temp} \PY{o}{=} \PY{n}{train\PYZus{}test\PYZus{}split}\PY{p}{(}
    \PY{n}{X}\PY{p}{,} \PY{n}{y}\PY{p}{,} \PY{n}{test\PYZus{}size}\PY{o}{=}\PY{l+m+mf}{0.4}\PY{p}{,} \PY{n}{random\PYZus{}state}\PY{o}{=}\PY{l+m+mi}{42}\PY{p}{,} \PY{n}{stratify}\PY{o}{=}\PY{n}{y}
\PY{p}{)}
\PY{n}{X\PYZus{}val}\PY{p}{,} \PY{n}{X\PYZus{}test}\PY{p}{,} \PY{n}{y\PYZus{}val}\PY{p}{,} \PY{n}{y\PYZus{}test} \PY{o}{=} \PY{n}{train\PYZus{}test\PYZus{}split}\PY{p}{(}
    \PY{n}{X\PYZus{}temp}\PY{p}{,} \PY{n}{y\PYZus{}temp}\PY{p}{,} \PY{n}{test\PYZus{}size}\PY{o}{=}\PY{l+m+mf}{0.5}\PY{p}{,} \PY{n}{random\PYZus{}state}\PY{o}{=}\PY{l+m+mi}{42}\PY{p}{,} \PY{n}{stratify}\PY{o}{=}\PY{n}{y\PYZus{}temp}
\PY{p}{)}

\PY{n+nb}{print}\PY{p}{(}\PY{l+s+sa}{f}\PY{l+s+s2}{\PYZdq{}}\PY{l+s+se}{\PYZbs{}n}\PY{l+s+s2}{Divisao: Treino=}\PY{l+s+si}{\PYZob{}}\PY{n+nb}{len}\PY{p}{(}\PY{n}{X\PYZus{}train}\PY{p}{)}\PY{l+s+si}{\PYZcb{}}\PY{l+s+s2}{, Validacao=}\PY{l+s+si}{\PYZob{}}\PY{n+nb}{len}\PY{p}{(}\PY{n}{X\PYZus{}val}\PY{p}{)}\PY{l+s+si}{\PYZcb{}}\PY{l+s+s2}{, Teste=}\PY{l+s+si}{\PYZob{}}\PY{n+nb}{len}\PY{p}{(}\PY{n}{X\PYZus{}test}\PY{p}{)}\PY{l+s+si}{\PYZcb{}}\PY{l+s+s2}{\PYZdq{}}\PY{p}{)}
\end{Verbatim}
\end{tcolorbox}

    \begin{Verbatim}[commandchars=\\\{\}]
Dimensoes: (569, 31)
Tipos de dados:
 float64    30
int64       1
Name: count, dtype: int64
Valores ausentes: 0
Linhas duplicadas: 0
Quantidade por classe:
 target
1    357
0    212
Name: count, dtype: int64

Divisao: Treino=341, Validacao=114, Teste=114
    \end{Verbatim}

    Resposta:

\begin{itemize}
\tightlist
\item
  Dimensoes: 569 linhas e 31 colunas (30 atributos continuos float64 e a
  variavel alvo inteira int64).
\item
  Valores ausentes: 0.
\item
  Linhas duplicadas: 0.
\item
  Quantidade por classe: 212 malignos (0) e 357 benignos (1).
\item
  Tratamento necessario: nenhum tratamento adicional e necessario, pois
  nao existem dados faltantes e os algoritmos baseados em arvores nao
  exigem normalizacao ou padronizacao de escala.
\item
  Preservar a proporcao das classes: a estratificacao garante que
  treino, validacao e teste tenham a mesma distribuicao original (cerca
  de 37.3\% malignos e 62.7\% benignos), evitando que subconjuntos
  fiquem com poucas amostras da classe minoritaria e gerem avaliacoes
  enviesadas.
\end{itemize}

    \subsection{Questão 2 --- Random
Forest}\label{questuxe3o-2-random-forest}

\begin{itemize}
\item
  Treine um
  \texttt{RandomForestClassifier(n\_estimators=100,\ random\_state=42)}.
  Na validação, apresente matriz de confusão com ordem
  \texttt{{[}0,\ 1{]}}, accuracy, precision, recall e F1 para
  \textbf{maligno} (\texttt{pos\_label=0}).
\item
  Teste pelo menos quatro combinações de \texttt{n\_estimators} (por
  exemplo, 50 e 200) e \texttt{max\_depth} (por exemplo, 3 e
  \texttt{None}). Compare os resultados de validação e explique o efeito
  dos parâmetros.
\end{itemize}

    \begin{tcolorbox}[breakable, size=fbox, boxrule=1pt, pad at break*=1mm,colback=cellbackground, colframe=cellborder]
\prompt{In}{incolor}{2}{\boxspacing}
\begin{Verbatim}[commandchars=\\\{\}]
\PY{k+kn}{from}\PY{+w}{ }\PY{n+nn}{sklearn}\PY{n+nn}{.}\PY{n+nn}{ensemble}\PY{+w}{ }\PY{k+kn}{import} \PY{n}{RandomForestClassifier}
\PY{k+kn}{from}\PY{+w}{ }\PY{n+nn}{sklearn}\PY{n+nn}{.}\PY{n+nn}{metrics}\PY{+w}{ }\PY{k+kn}{import} \PY{n}{accuracy\PYZus{}score}\PY{p}{,} \PY{n}{precision\PYZus{}score}\PY{p}{,} \PY{n}{recall\PYZus{}score}\PY{p}{,} \PY{n}{f1\PYZus{}score}\PY{p}{,} \PY{n}{confusion\PYZus{}matrix}

\PY{n}{rf} \PY{o}{=} \PY{n}{RandomForestClassifier}\PY{p}{(}\PY{n}{n\PYZus{}estimators}\PY{o}{=}\PY{l+m+mi}{100}\PY{p}{,} \PY{n}{random\PYZus{}state}\PY{o}{=}\PY{l+m+mi}{42}\PY{p}{)}
\PY{n}{rf}\PY{o}{.}\PY{n}{fit}\PY{p}{(}\PY{n}{X\PYZus{}train}\PY{p}{,} \PY{n}{y\PYZus{}train}\PY{p}{)}
\PY{n}{y\PYZus{}pred\PYZus{}val} \PY{o}{=} \PY{n}{rf}\PY{o}{.}\PY{n}{predict}\PY{p}{(}\PY{n}{X\PYZus{}val}\PY{p}{)}

\PY{n+nb}{print}\PY{p}{(}\PY{l+s+s2}{\PYZdq{}}\PY{l+s+s2}{Matriz de Confusao (ordem [0, 1]):}\PY{l+s+se}{\PYZbs{}n}\PY{l+s+s2}{\PYZdq{}}\PY{p}{,} \PY{n}{confusion\PYZus{}matrix}\PY{p}{(}\PY{n}{y\PYZus{}val}\PY{p}{,} \PY{n}{y\PYZus{}pred\PYZus{}val}\PY{p}{,} \PY{n}{labels}\PY{o}{=}\PY{p}{[}\PY{l+m+mi}{0}\PY{p}{,} \PY{l+m+mi}{1}\PY{p}{]}\PY{p}{)}\PY{p}{)}
\PY{n+nb}{print}\PY{p}{(}\PY{l+s+sa}{f}\PY{l+s+s2}{\PYZdq{}}\PY{l+s+s2}{Accuracy: }\PY{l+s+si}{\PYZob{}}\PY{n}{accuracy\PYZus{}score}\PY{p}{(}\PY{n}{y\PYZus{}val}\PY{p}{,}\PY{+w}{ }\PY{n}{y\PYZus{}pred\PYZus{}val}\PY{p}{)}\PY{l+s+si}{:}\PY{l+s+s2}{.4f}\PY{l+s+si}{\PYZcb{}}\PY{l+s+s2}{\PYZdq{}}\PY{p}{)}
\PY{n+nb}{print}\PY{p}{(}\PY{l+s+sa}{f}\PY{l+s+s2}{\PYZdq{}}\PY{l+s+s2}{Precision (0): }\PY{l+s+si}{\PYZob{}}\PY{n}{precision\PYZus{}score}\PY{p}{(}\PY{n}{y\PYZus{}val}\PY{p}{,}\PY{+w}{ }\PY{n}{y\PYZus{}pred\PYZus{}val}\PY{p}{,}\PY{+w}{ }\PY{n}{pos\PYZus{}label}\PY{o}{=}\PY{l+m+mi}{0}\PY{p}{)}\PY{l+s+si}{:}\PY{l+s+s2}{.4f}\PY{l+s+si}{\PYZcb{}}\PY{l+s+s2}{\PYZdq{}}\PY{p}{)}
\PY{n+nb}{print}\PY{p}{(}\PY{l+s+sa}{f}\PY{l+s+s2}{\PYZdq{}}\PY{l+s+s2}{Recall (0): }\PY{l+s+si}{\PYZob{}}\PY{n}{recall\PYZus{}score}\PY{p}{(}\PY{n}{y\PYZus{}val}\PY{p}{,}\PY{+w}{ }\PY{n}{y\PYZus{}pred\PYZus{}val}\PY{p}{,}\PY{+w}{ }\PY{n}{pos\PYZus{}label}\PY{o}{=}\PY{l+m+mi}{0}\PY{p}{)}\PY{l+s+si}{:}\PY{l+s+s2}{.4f}\PY{l+s+si}{\PYZcb{}}\PY{l+s+s2}{\PYZdq{}}\PY{p}{)}
\PY{n+nb}{print}\PY{p}{(}\PY{l+s+sa}{f}\PY{l+s+s2}{\PYZdq{}}\PY{l+s+s2}{F1 (0): }\PY{l+s+si}{\PYZob{}}\PY{n}{f1\PYZus{}score}\PY{p}{(}\PY{n}{y\PYZus{}val}\PY{p}{,}\PY{+w}{ }\PY{n}{y\PYZus{}pred\PYZus{}val}\PY{p}{,}\PY{+w}{ }\PY{n}{pos\PYZus{}label}\PY{o}{=}\PY{l+m+mi}{0}\PY{p}{)}\PY{l+s+si}{:}\PY{l+s+s2}{.4f}\PY{l+s+si}{\PYZcb{}}\PY{l+s+s2}{\PYZdq{}}\PY{p}{)}

\PY{n}{combos} \PY{o}{=} \PY{p}{[}\PY{p}{(}\PY{l+m+mi}{50}\PY{p}{,} \PY{l+m+mi}{3}\PY{p}{)}\PY{p}{,} \PY{p}{(}\PY{l+m+mi}{50}\PY{p}{,} \PY{k+kc}{None}\PY{p}{)}\PY{p}{,} \PY{p}{(}\PY{l+m+mi}{200}\PY{p}{,} \PY{l+m+mi}{3}\PY{p}{)}\PY{p}{,} \PY{p}{(}\PY{l+m+mi}{200}\PY{p}{,} \PY{k+kc}{None}\PY{p}{)}\PY{p}{]}
\PY{n}{rf\PYZus{}results} \PY{o}{=} \PY{p}{[}\PY{p}{]}

\PY{k}{for} \PY{n}{n\PYZus{}est}\PY{p}{,} \PY{n}{depth} \PY{o+ow}{in} \PY{n}{combos}\PY{p}{:}
    \PY{n}{m} \PY{o}{=} \PY{n}{RandomForestClassifier}\PY{p}{(}\PY{n}{n\PYZus{}estimators}\PY{o}{=}\PY{n}{n\PYZus{}est}\PY{p}{,} \PY{n}{max\PYZus{}depth}\PY{o}{=}\PY{n}{depth}\PY{p}{,} \PY{n}{random\PYZus{}state}\PY{o}{=}\PY{l+m+mi}{42}\PY{p}{)}
    \PY{n}{m}\PY{o}{.}\PY{n}{fit}\PY{p}{(}\PY{n}{X\PYZus{}train}\PY{p}{,} \PY{n}{y\PYZus{}train}\PY{p}{)}
    \PY{n}{p} \PY{o}{=} \PY{n}{m}\PY{o}{.}\PY{n}{predict}\PY{p}{(}\PY{n}{X\PYZus{}val}\PY{p}{)}
    \PY{n}{rf\PYZus{}results}\PY{o}{.}\PY{n}{append}\PY{p}{(}\PY{p}{\PYZob{}}
        \PY{l+s+s1}{\PYZsq{}}\PY{l+s+s1}{n\PYZus{}estimators}\PY{l+s+s1}{\PYZsq{}}\PY{p}{:} \PY{n}{n\PYZus{}est}\PY{p}{,}
        \PY{l+s+s1}{\PYZsq{}}\PY{l+s+s1}{max\PYZus{}depth}\PY{l+s+s1}{\PYZsq{}}\PY{p}{:} \PY{n+nb}{str}\PY{p}{(}\PY{n}{depth}\PY{p}{)}\PY{p}{,}
        \PY{l+s+s1}{\PYZsq{}}\PY{l+s+s1}{Accuracy}\PY{l+s+s1}{\PYZsq{}}\PY{p}{:} \PY{n}{accuracy\PYZus{}score}\PY{p}{(}\PY{n}{y\PYZus{}val}\PY{p}{,} \PY{n}{p}\PY{p}{)}\PY{p}{,}
        \PY{l+s+s1}{\PYZsq{}}\PY{l+s+s1}{Precision}\PY{l+s+s1}{\PYZsq{}}\PY{p}{:} \PY{n}{precision\PYZus{}score}\PY{p}{(}\PY{n}{y\PYZus{}val}\PY{p}{,} \PY{n}{p}\PY{p}{,} \PY{n}{pos\PYZus{}label}\PY{o}{=}\PY{l+m+mi}{0}\PY{p}{)}\PY{p}{,}
        \PY{l+s+s1}{\PYZsq{}}\PY{l+s+s1}{Recall}\PY{l+s+s1}{\PYZsq{}}\PY{p}{:} \PY{n}{recall\PYZus{}score}\PY{p}{(}\PY{n}{y\PYZus{}val}\PY{p}{,} \PY{n}{p}\PY{p}{,} \PY{n}{pos\PYZus{}label}\PY{o}{=}\PY{l+m+mi}{0}\PY{p}{)}\PY{p}{,}
        \PY{l+s+s1}{\PYZsq{}}\PY{l+s+s1}{F1}\PY{l+s+s1}{\PYZsq{}}\PY{p}{:} \PY{n}{f1\PYZus{}score}\PY{p}{(}\PY{n}{y\PYZus{}val}\PY{p}{,} \PY{n}{p}\PY{p}{,} \PY{n}{pos\PYZus{}label}\PY{o}{=}\PY{l+m+mi}{0}\PY{p}{)}
    \PY{p}{\PYZcb{}}\PY{p}{)}

\PY{n}{pd}\PY{o}{.}\PY{n}{DataFrame}\PY{p}{(}\PY{n}{rf\PYZus{}results}\PY{p}{)}
\end{Verbatim}
\end{tcolorbox}

    \begin{Verbatim}[commandchars=\\\{\}]
Matriz de Confusao (ordem [0, 1]):
 [[39  3]
 [ 1 71]]
Accuracy: 0.9649
Precision (0): 0.9750
Recall (0): 0.9286
F1 (0): 0.9512
    \end{Verbatim}

            \begin{tcolorbox}[breakable, size=fbox, boxrule=.5pt, pad at break*=1mm, opacityfill=0]
\prompt{Out}{outcolor}{2}{\boxspacing}
\begin{Verbatim}[commandchars=\\\{\}]
   n\_estimators max\_depth  Accuracy  Precision    Recall       F1
0            50         3  0.964912      0.975  0.928571  0.95122
1            50      None  0.964912      0.975  0.928571  0.95122
2           200         3  0.964912      0.975  0.928571  0.95122
3           200      None  0.964912      0.975  0.928571  0.95122
\end{Verbatim}
\end{tcolorbox}
        
    Resposta:

\begin{itemize}
\tightlist
\item
  Resultados na validacao do modelo inicial: Matriz de confusao
  {[}{[}39, 3{]}, {[}1, 71{]}{]}, Acuracia = 0.9649, Precisao = 0.9750,
  Recall = 0.9286 e F1 = 0.9512.
\item
  Efeito dos parametros:

  \begin{itemize}
  \tightlist
  \item
    n\_estimators: numero de arvores na floresta. Aumentar esse valor
    reduz a variancia e estabiliza o modelo, sem gerar overfitting, mas
    com maior custo computacional.
  \item
    max\_depth: profundidade maxima das arvores. Controla a complexidade
    individual das arvores. max\_depth=3 limita o crescimento evitando
    memorizacao, enquanto None permite crescimento total.
  \item
    No conjunto de validacao, todas as combinacoes obtiveram desempenho
    identico (F1 = 0.9512), mostrando que max\_depth=3 ja e suficiente
    para obter a separacao dos dados com menor complexidade.
  \end{itemize}
\end{itemize}

    \subsection{Questão 3 --- AdaBoost}\label{questuxe3o-3-adaboost}

\begin{itemize}
\item
  Treine um
  \texttt{AdaBoostClassifier(n\_estimators=50,\ learning\_rate=1.0,\ random\_state=42)}.
\item
  obs: \texttt{learning\_rate}: Peso atribuído a cada classificador em
  cada iteração de \emph{boosting}. Uma taxa de aprendizado mais alta
  aumenta a contribuição de cada classificador. Existe um
  \emph{trade-off} entre os parâmetros \texttt{learning\_rate} e
  \texttt{n\_estimators}.
\item
  Calcule na validação as mesmas métricas e a matriz de confusão da
  questão 2. Teste pelo menos quatro combinações de
  \texttt{n\_estimators} (por exemplo, 50 e 100) e
  \texttt{learning\_rate} (por exemplo, 0.1 e 1.0).
\item
  Como os parâmetros alteraram o desempenho? Compare o mecanismo de
  boosting do AdaBoost com o bagging usado no Random Forest.
\end{itemize}

    \begin{tcolorbox}[breakable, size=fbox, boxrule=1pt, pad at break*=1mm,colback=cellbackground, colframe=cellborder]
\prompt{In}{incolor}{3}{\boxspacing}
\begin{Verbatim}[commandchars=\\\{\}]
\PY{k+kn}{from}\PY{+w}{ }\PY{n+nn}{sklearn}\PY{n+nn}{.}\PY{n+nn}{ensemble}\PY{+w}{ }\PY{k+kn}{import} \PY{n}{AdaBoostClassifier}

\PY{n}{ada} \PY{o}{=} \PY{n}{AdaBoostClassifier}\PY{p}{(}\PY{n}{n\PYZus{}estimators}\PY{o}{=}\PY{l+m+mi}{50}\PY{p}{,} \PY{n}{learning\PYZus{}rate}\PY{o}{=}\PY{l+m+mf}{1.0}\PY{p}{,} \PY{n}{random\PYZus{}state}\PY{o}{=}\PY{l+m+mi}{42}\PY{p}{)}
\PY{n}{ada}\PY{o}{.}\PY{n}{fit}\PY{p}{(}\PY{n}{X\PYZus{}train}\PY{p}{,} \PY{n}{y\PYZus{}train}\PY{p}{)}
\PY{n}{y\PYZus{}pred\PYZus{}val\PYZus{}ada} \PY{o}{=} \PY{n}{ada}\PY{o}{.}\PY{n}{predict}\PY{p}{(}\PY{n}{X\PYZus{}val}\PY{p}{)}

\PY{n+nb}{print}\PY{p}{(}\PY{l+s+s2}{\PYZdq{}}\PY{l+s+s2}{Matriz de Confusao (ordem [0, 1]):}\PY{l+s+se}{\PYZbs{}n}\PY{l+s+s2}{\PYZdq{}}\PY{p}{,} \PY{n}{confusion\PYZus{}matrix}\PY{p}{(}\PY{n}{y\PYZus{}val}\PY{p}{,} \PY{n}{y\PYZus{}pred\PYZus{}val\PYZus{}ada}\PY{p}{,} \PY{n}{labels}\PY{o}{=}\PY{p}{[}\PY{l+m+mi}{0}\PY{p}{,} \PY{l+m+mi}{1}\PY{p}{]}\PY{p}{)}\PY{p}{)}
\PY{n+nb}{print}\PY{p}{(}\PY{l+s+sa}{f}\PY{l+s+s2}{\PYZdq{}}\PY{l+s+s2}{Accuracy: }\PY{l+s+si}{\PYZob{}}\PY{n}{accuracy\PYZus{}score}\PY{p}{(}\PY{n}{y\PYZus{}val}\PY{p}{,}\PY{+w}{ }\PY{n}{y\PYZus{}pred\PYZus{}val\PYZus{}ada}\PY{p}{)}\PY{l+s+si}{:}\PY{l+s+s2}{.4f}\PY{l+s+si}{\PYZcb{}}\PY{l+s+s2}{\PYZdq{}}\PY{p}{)}
\PY{n+nb}{print}\PY{p}{(}\PY{l+s+sa}{f}\PY{l+s+s2}{\PYZdq{}}\PY{l+s+s2}{Precision (0): }\PY{l+s+si}{\PYZob{}}\PY{n}{precision\PYZus{}score}\PY{p}{(}\PY{n}{y\PYZus{}val}\PY{p}{,}\PY{+w}{ }\PY{n}{y\PYZus{}pred\PYZus{}val\PYZus{}ada}\PY{p}{,}\PY{+w}{ }\PY{n}{pos\PYZus{}label}\PY{o}{=}\PY{l+m+mi}{0}\PY{p}{)}\PY{l+s+si}{:}\PY{l+s+s2}{.4f}\PY{l+s+si}{\PYZcb{}}\PY{l+s+s2}{\PYZdq{}}\PY{p}{)}
\PY{n+nb}{print}\PY{p}{(}\PY{l+s+sa}{f}\PY{l+s+s2}{\PYZdq{}}\PY{l+s+s2}{Recall (0): }\PY{l+s+si}{\PYZob{}}\PY{n}{recall\PYZus{}score}\PY{p}{(}\PY{n}{y\PYZus{}val}\PY{p}{,}\PY{+w}{ }\PY{n}{y\PYZus{}pred\PYZus{}val\PYZus{}ada}\PY{p}{,}\PY{+w}{ }\PY{n}{pos\PYZus{}label}\PY{o}{=}\PY{l+m+mi}{0}\PY{p}{)}\PY{l+s+si}{:}\PY{l+s+s2}{.4f}\PY{l+s+si}{\PYZcb{}}\PY{l+s+s2}{\PYZdq{}}\PY{p}{)}
\PY{n+nb}{print}\PY{p}{(}\PY{l+s+sa}{f}\PY{l+s+s2}{\PYZdq{}}\PY{l+s+s2}{F1 (0): }\PY{l+s+si}{\PYZob{}}\PY{n}{f1\PYZus{}score}\PY{p}{(}\PY{n}{y\PYZus{}val}\PY{p}{,}\PY{+w}{ }\PY{n}{y\PYZus{}pred\PYZus{}val\PYZus{}ada}\PY{p}{,}\PY{+w}{ }\PY{n}{pos\PYZus{}label}\PY{o}{=}\PY{l+m+mi}{0}\PY{p}{)}\PY{l+s+si}{:}\PY{l+s+s2}{.4f}\PY{l+s+si}{\PYZcb{}}\PY{l+s+s2}{\PYZdq{}}\PY{p}{)}

\PY{n}{ada\PYZus{}combos} \PY{o}{=} \PY{p}{[}\PY{p}{(}\PY{l+m+mi}{50}\PY{p}{,} \PY{l+m+mf}{0.1}\PY{p}{)}\PY{p}{,} \PY{p}{(}\PY{l+m+mi}{50}\PY{p}{,} \PY{l+m+mf}{1.0}\PY{p}{)}\PY{p}{,} \PY{p}{(}\PY{l+m+mi}{100}\PY{p}{,} \PY{l+m+mf}{0.1}\PY{p}{)}\PY{p}{,} \PY{p}{(}\PY{l+m+mi}{100}\PY{p}{,} \PY{l+m+mf}{1.0}\PY{p}{)}\PY{p}{]}
\PY{n}{ada\PYZus{}results} \PY{o}{=} \PY{p}{[}\PY{p}{]}

\PY{k}{for} \PY{n}{n\PYZus{}est}\PY{p}{,} \PY{n}{lr} \PY{o+ow}{in} \PY{n}{ada\PYZus{}combos}\PY{p}{:}
    \PY{n}{m} \PY{o}{=} \PY{n}{AdaBoostClassifier}\PY{p}{(}\PY{n}{n\PYZus{}estimators}\PY{o}{=}\PY{n}{n\PYZus{}est}\PY{p}{,} \PY{n}{learning\PYZus{}rate}\PY{o}{=}\PY{n}{lr}\PY{p}{,} \PY{n}{random\PYZus{}state}\PY{o}{=}\PY{l+m+mi}{42}\PY{p}{)}
    \PY{n}{m}\PY{o}{.}\PY{n}{fit}\PY{p}{(}\PY{n}{X\PYZus{}train}\PY{p}{,} \PY{n}{y\PYZus{}train}\PY{p}{)}
    \PY{n}{p} \PY{o}{=} \PY{n}{m}\PY{o}{.}\PY{n}{predict}\PY{p}{(}\PY{n}{X\PYZus{}val}\PY{p}{)}
    \PY{n}{ada\PYZus{}results}\PY{o}{.}\PY{n}{append}\PY{p}{(}\PY{p}{\PYZob{}}
        \PY{l+s+s1}{\PYZsq{}}\PY{l+s+s1}{n\PYZus{}estimators}\PY{l+s+s1}{\PYZsq{}}\PY{p}{:} \PY{n}{n\PYZus{}est}\PY{p}{,}
        \PY{l+s+s1}{\PYZsq{}}\PY{l+s+s1}{learning\PYZus{}rate}\PY{l+s+s1}{\PYZsq{}}\PY{p}{:} \PY{n}{lr}\PY{p}{,}
        \PY{l+s+s1}{\PYZsq{}}\PY{l+s+s1}{Accuracy}\PY{l+s+s1}{\PYZsq{}}\PY{p}{:} \PY{n}{accuracy\PYZus{}score}\PY{p}{(}\PY{n}{y\PYZus{}val}\PY{p}{,} \PY{n}{p}\PY{p}{)}\PY{p}{,}
        \PY{l+s+s1}{\PYZsq{}}\PY{l+s+s1}{Precision}\PY{l+s+s1}{\PYZsq{}}\PY{p}{:} \PY{n}{precision\PYZus{}score}\PY{p}{(}\PY{n}{y\PYZus{}val}\PY{p}{,} \PY{n}{p}\PY{p}{,} \PY{n}{pos\PYZus{}label}\PY{o}{=}\PY{l+m+mi}{0}\PY{p}{)}\PY{p}{,}
        \PY{l+s+s1}{\PYZsq{}}\PY{l+s+s1}{Recall}\PY{l+s+s1}{\PYZsq{}}\PY{p}{:} \PY{n}{recall\PYZus{}score}\PY{p}{(}\PY{n}{y\PYZus{}val}\PY{p}{,} \PY{n}{p}\PY{p}{,} \PY{n}{pos\PYZus{}label}\PY{o}{=}\PY{l+m+mi}{0}\PY{p}{)}\PY{p}{,}
        \PY{l+s+s1}{\PYZsq{}}\PY{l+s+s1}{F1}\PY{l+s+s1}{\PYZsq{}}\PY{p}{:} \PY{n}{f1\PYZus{}score}\PY{p}{(}\PY{n}{y\PYZus{}val}\PY{p}{,} \PY{n}{p}\PY{p}{,} \PY{n}{pos\PYZus{}label}\PY{o}{=}\PY{l+m+mi}{0}\PY{p}{)}
    \PY{p}{\PYZcb{}}\PY{p}{)}

\PY{n}{pd}\PY{o}{.}\PY{n}{DataFrame}\PY{p}{(}\PY{n}{ada\PYZus{}results}\PY{p}{)}
\end{Verbatim}
\end{tcolorbox}

    \begin{Verbatim}[commandchars=\\\{\}]
Matriz de Confusao (ordem [0, 1]):
 [[39  3]
 [ 1 71]]
Accuracy: 0.9649
Precision (0): 0.9750
Recall (0): 0.9286
F1 (0): 0.9512
    \end{Verbatim}

            \begin{tcolorbox}[breakable, size=fbox, boxrule=.5pt, pad at break*=1mm, opacityfill=0]
\prompt{Out}{outcolor}{3}{\boxspacing}
\begin{Verbatim}[commandchars=\\\{\}]
   n\_estimators  learning\_rate  Accuracy  Precision    Recall        F1
0            50            0.1  0.956140    0.95122  0.928571  0.939759
1            50            1.0  0.964912    0.97500  0.928571  0.951220
2           100            0.1  0.964912    0.97500  0.928571  0.951220
3           100            1.0  0.964912    0.97500  0.928571  0.951220
\end{Verbatim}
\end{tcolorbox}
        
    Resposta:

\begin{itemize}
\tightlist
\item
  Resultados na validacao do AdaBoost inicial: Matriz de confusao
  {[}{[}39, 3{]}, {[}1, 71{]}{]}, Acuracia = 0.9649, Precisao = 0.9750,
  Recall = 0.9286 e F1 = 0.9512.
\item
  Efeito dos parametros: com learning\_rate=0.1 e n\_estimators=50 o
  modelo teve aprendizado mais lento e F1 menor (0.9398); ao aumentar
  para n\_estimators=100 com learning\_rate=0.1 ou usar
  learning\_rate=1.0 com 50 estimadores, o F1 alcancou 0.9512,
  demonstrando o trade-off entre taxa de aprendizado e numero de
  estimadores.
\item
  Comparacao Boosting vs Bagging:

  \begin{itemize}
  \tightlist
  \item
    Bagging (Random Forest): treina arvores completas em paralelo e de
    forma independente sobre amostras bootstrap, visando diminuir a
    variancia.
  \item
    Boosting (AdaBoost): treina estimadores fracos em sequencia, dando
    mais peso as amostras erradas nas iteracoes anteriores, visando
    diminuir o vies.
  \end{itemize}
\end{itemize}

    \subsection{Questão 4 --- Analise os erros e escolha
configurações}\label{questuxe3o-4-analise-os-erros-e-escolha-configurauxe7uxf5es}

\begin{itemize}
\tightlist
\item
  Construa uma tabela de validação com os modelos iniciais e as
  configurações testadas, contendo accuracy, precision, recall, F1 para
  a classe maligna e número de \textbf{malignos previstos como
  benignos}. Explique por que accuracy sozinha pode ocultar um erro
  relevante.
\item
  Escolha uma configuração final de cada algoritmo com um critério
  explícito; em caso de empate, explique como desempata. Não use o
  conjunto de teste para esta escolha. (Aqui você pode usar o
  \texttt{RandomizedSearchCV} do scikit-learn para busca dos melhores
  hiper-parâmetros, que testa combinações aleatórias de hiperparâmetros
  dentro de um intervalo definido.)
\end{itemize}

    \begin{tcolorbox}[breakable, size=fbox, boxrule=1pt, pad at break*=1mm,colback=cellbackground, colframe=cellborder]
\prompt{In}{incolor}{4}{\boxspacing}
\begin{Verbatim}[commandchars=\\\{\}]
\PY{n}{models} \PY{o}{=} \PY{p}{[}
    \PY{p}{(}\PY{l+s+s1}{\PYZsq{}}\PY{l+s+s1}{RF n=100, d=None}\PY{l+s+s1}{\PYZsq{}}\PY{p}{,} \PY{n}{RandomForestClassifier}\PY{p}{(}\PY{n}{n\PYZus{}estimators}\PY{o}{=}\PY{l+m+mi}{100}\PY{p}{,} \PY{n}{random\PYZus{}state}\PY{o}{=}\PY{l+m+mi}{42}\PY{p}{)}\PY{p}{)}\PY{p}{,}
    \PY{p}{(}\PY{l+s+s1}{\PYZsq{}}\PY{l+s+s1}{RF n=50, d=3}\PY{l+s+s1}{\PYZsq{}}\PY{p}{,} \PY{n}{RandomForestClassifier}\PY{p}{(}\PY{n}{n\PYZus{}estimators}\PY{o}{=}\PY{l+m+mi}{50}\PY{p}{,} \PY{n}{max\PYZus{}depth}\PY{o}{=}\PY{l+m+mi}{3}\PY{p}{,} \PY{n}{random\PYZus{}state}\PY{o}{=}\PY{l+m+mi}{42}\PY{p}{)}\PY{p}{)}\PY{p}{,}
    \PY{p}{(}\PY{l+s+s1}{\PYZsq{}}\PY{l+s+s1}{RF n=50, d=None}\PY{l+s+s1}{\PYZsq{}}\PY{p}{,} \PY{n}{RandomForestClassifier}\PY{p}{(}\PY{n}{n\PYZus{}estimators}\PY{o}{=}\PY{l+m+mi}{50}\PY{p}{,} \PY{n}{max\PYZus{}depth}\PY{o}{=}\PY{k+kc}{None}\PY{p}{,} \PY{n}{random\PYZus{}state}\PY{o}{=}\PY{l+m+mi}{42}\PY{p}{)}\PY{p}{)}\PY{p}{,}
    \PY{p}{(}\PY{l+s+s1}{\PYZsq{}}\PY{l+s+s1}{RF n=200, d=3}\PY{l+s+s1}{\PYZsq{}}\PY{p}{,} \PY{n}{RandomForestClassifier}\PY{p}{(}\PY{n}{n\PYZus{}estimators}\PY{o}{=}\PY{l+m+mi}{200}\PY{p}{,} \PY{n}{max\PYZus{}depth}\PY{o}{=}\PY{l+m+mi}{3}\PY{p}{,} \PY{n}{random\PYZus{}state}\PY{o}{=}\PY{l+m+mi}{42}\PY{p}{)}\PY{p}{)}\PY{p}{,}
    \PY{p}{(}\PY{l+s+s1}{\PYZsq{}}\PY{l+s+s1}{RF n=200, d=None}\PY{l+s+s1}{\PYZsq{}}\PY{p}{,} \PY{n}{RandomForestClassifier}\PY{p}{(}\PY{n}{n\PYZus{}estimators}\PY{o}{=}\PY{l+m+mi}{200}\PY{p}{,} \PY{n}{max\PYZus{}depth}\PY{o}{=}\PY{k+kc}{None}\PY{p}{,} \PY{n}{random\PYZus{}state}\PY{o}{=}\PY{l+m+mi}{42}\PY{p}{)}\PY{p}{)}\PY{p}{,}
    \PY{p}{(}\PY{l+s+s1}{\PYZsq{}}\PY{l+s+s1}{Ada n=50, lr=1.0}\PY{l+s+s1}{\PYZsq{}}\PY{p}{,} \PY{n}{AdaBoostClassifier}\PY{p}{(}\PY{n}{n\PYZus{}estimators}\PY{o}{=}\PY{l+m+mi}{50}\PY{p}{,} \PY{n}{learning\PYZus{}rate}\PY{o}{=}\PY{l+m+mf}{1.0}\PY{p}{,} \PY{n}{random\PYZus{}state}\PY{o}{=}\PY{l+m+mi}{42}\PY{p}{)}\PY{p}{)}\PY{p}{,}
    \PY{p}{(}\PY{l+s+s1}{\PYZsq{}}\PY{l+s+s1}{Ada n=50, lr=0.1}\PY{l+s+s1}{\PYZsq{}}\PY{p}{,} \PY{n}{AdaBoostClassifier}\PY{p}{(}\PY{n}{n\PYZus{}estimators}\PY{o}{=}\PY{l+m+mi}{50}\PY{p}{,} \PY{n}{learning\PYZus{}rate}\PY{o}{=}\PY{l+m+mf}{0.1}\PY{p}{,} \PY{n}{random\PYZus{}state}\PY{o}{=}\PY{l+m+mi}{42}\PY{p}{)}\PY{p}{)}\PY{p}{,}
    \PY{p}{(}\PY{l+s+s1}{\PYZsq{}}\PY{l+s+s1}{Ada n=100, lr=0.1}\PY{l+s+s1}{\PYZsq{}}\PY{p}{,} \PY{n}{AdaBoostClassifier}\PY{p}{(}\PY{n}{n\PYZus{}estimators}\PY{o}{=}\PY{l+m+mi}{100}\PY{p}{,} \PY{n}{learning\PYZus{}rate}\PY{o}{=}\PY{l+m+mf}{0.1}\PY{p}{,} \PY{n}{random\PYZus{}state}\PY{o}{=}\PY{l+m+mi}{42}\PY{p}{)}\PY{p}{)}\PY{p}{,}
    \PY{p}{(}\PY{l+s+s1}{\PYZsq{}}\PY{l+s+s1}{Ada n=100, lr=1.0}\PY{l+s+s1}{\PYZsq{}}\PY{p}{,} \PY{n}{AdaBoostClassifier}\PY{p}{(}\PY{n}{n\PYZus{}estimators}\PY{o}{=}\PY{l+m+mi}{100}\PY{p}{,} \PY{n}{learning\PYZus{}rate}\PY{o}{=}\PY{l+m+mf}{1.0}\PY{p}{,} \PY{n}{random\PYZus{}state}\PY{o}{=}\PY{l+m+mi}{42}\PY{p}{)}\PY{p}{)}
\PY{p}{]}

\PY{n}{summary} \PY{o}{=} \PY{p}{[}\PY{p}{]}
\PY{k}{for} \PY{n}{name}\PY{p}{,} \PY{n}{m} \PY{o+ow}{in} \PY{n}{models}\PY{p}{:}
    \PY{n}{m}\PY{o}{.}\PY{n}{fit}\PY{p}{(}\PY{n}{X\PYZus{}train}\PY{p}{,} \PY{n}{y\PYZus{}train}\PY{p}{)}
    \PY{n}{p} \PY{o}{=} \PY{n}{m}\PY{o}{.}\PY{n}{predict}\PY{p}{(}\PY{n}{X\PYZus{}val}\PY{p}{)}
    \PY{n}{cm} \PY{o}{=} \PY{n}{confusion\PYZus{}matrix}\PY{p}{(}\PY{n}{y\PYZus{}val}\PY{p}{,} \PY{n}{p}\PY{p}{,} \PY{n}{labels}\PY{o}{=}\PY{p}{[}\PY{l+m+mi}{0}\PY{p}{,} \PY{l+m+mi}{1}\PY{p}{]}\PY{p}{)}
    \PY{n}{summary}\PY{o}{.}\PY{n}{append}\PY{p}{(}\PY{p}{\PYZob{}}
        \PY{l+s+s1}{\PYZsq{}}\PY{l+s+s1}{Modelo}\PY{l+s+s1}{\PYZsq{}}\PY{p}{:} \PY{n}{name}\PY{p}{,}
        \PY{l+s+s1}{\PYZsq{}}\PY{l+s+s1}{Accuracy}\PY{l+s+s1}{\PYZsq{}}\PY{p}{:} \PY{n}{accuracy\PYZus{}score}\PY{p}{(}\PY{n}{y\PYZus{}val}\PY{p}{,} \PY{n}{p}\PY{p}{)}\PY{p}{,}
        \PY{l+s+s1}{\PYZsq{}}\PY{l+s+s1}{Precision (0)}\PY{l+s+s1}{\PYZsq{}}\PY{p}{:} \PY{n}{precision\PYZus{}score}\PY{p}{(}\PY{n}{y\PYZus{}val}\PY{p}{,} \PY{n}{p}\PY{p}{,} \PY{n}{pos\PYZus{}label}\PY{o}{=}\PY{l+m+mi}{0}\PY{p}{)}\PY{p}{,}
        \PY{l+s+s1}{\PYZsq{}}\PY{l+s+s1}{Recall (0)}\PY{l+s+s1}{\PYZsq{}}\PY{p}{:} \PY{n}{recall\PYZus{}score}\PY{p}{(}\PY{n}{y\PYZus{}val}\PY{p}{,} \PY{n}{p}\PY{p}{,} \PY{n}{pos\PYZus{}label}\PY{o}{=}\PY{l+m+mi}{0}\PY{p}{)}\PY{p}{,}
        \PY{l+s+s1}{\PYZsq{}}\PY{l+s+s1}{F1 (0)}\PY{l+s+s1}{\PYZsq{}}\PY{p}{:} \PY{n}{f1\PYZus{}score}\PY{p}{(}\PY{n}{y\PYZus{}val}\PY{p}{,} \PY{n}{p}\PY{p}{,} \PY{n}{pos\PYZus{}label}\PY{o}{=}\PY{l+m+mi}{0}\PY{p}{)}\PY{p}{,}
        \PY{l+s+s1}{\PYZsq{}}\PY{l+s+s1}{Malignos como benignos (FN)}\PY{l+s+s1}{\PYZsq{}}\PY{p}{:} \PY{n}{cm}\PY{p}{[}\PY{l+m+mi}{0}\PY{p}{,} \PY{l+m+mi}{1}\PY{p}{]}
    \PY{p}{\PYZcb{}}\PY{p}{)}

\PY{n}{pd}\PY{o}{.}\PY{n}{DataFrame}\PY{p}{(}\PY{n}{summary}\PY{p}{)}
\end{Verbatim}
\end{tcolorbox}

            \begin{tcolorbox}[breakable, size=fbox, boxrule=.5pt, pad at break*=1mm, opacityfill=0]
\prompt{Out}{outcolor}{4}{\boxspacing}
\begin{Verbatim}[commandchars=\\\{\}]
              Modelo  Accuracy  Precision (0)  Recall (0)    F1 (0)  \textbackslash{}
0   RF n=100, d=None  0.964912        0.97500    0.928571  0.951220
1       RF n=50, d=3  0.964912        0.97500    0.928571  0.951220
2    RF n=50, d=None  0.964912        0.97500    0.928571  0.951220
3      RF n=200, d=3  0.964912        0.97500    0.928571  0.951220
4   RF n=200, d=None  0.964912        0.97500    0.928571  0.951220
5   Ada n=50, lr=1.0  0.964912        0.97500    0.928571  0.951220
6   Ada n=50, lr=0.1  0.956140        0.95122    0.928571  0.939759
7  Ada n=100, lr=0.1  0.964912        0.97500    0.928571  0.951220
8  Ada n=100, lr=1.0  0.964912        0.97500    0.928571  0.951220

   Malignos como benignos (FN)
0                            3
1                            3
2                            3
3                            3
4                            3
5                            3
6                            3
7                            3
8                            3
\end{Verbatim}
\end{tcolorbox}
        
    Resposta:

\begin{itemize}
\tightlist
\item
  Por que acuracia sozinha pode ocultar erros: em contexto medico, falso
  negativo (maligno classificado como benigno) e o erro mais critico,
  pois atrasa o tratamento. A acuracia da o mesmo peso a todos os erros
  e pode ser alta mesmo se o modelo errar casos graves.
\item
  Criterio de escolha e desempate: prioriza-se o menor numero de falsos
  negativos (maior Recall da classe 0) e maior F1 da classe 0. Em caso
  de empate, escolhe-se o modelo mais simples (menor numero de
  estimadores e menor profundidade) para economizar recursos
  computacionais e evitar overfitting.
\item
  Modelos escolhidos:

  \begin{itemize}
  \tightlist
  \item
    Random Forest: n\_estimators=50, max\_depth=3 (obteve mesmo F1 de
    0.9512 e 3 falsos negativos com modelo mais simples).
  \item
    AdaBoost: n\_estimators=50, learning\_rate=1.0 (obteve F1 de 0.9512
    e 3 falsos negativos com menor numero de estimadores).
  \end{itemize}
\end{itemize}

    \subsection{Questão 5 --- Avaliação
final}\label{questuxe3o-5-avaliauxe7uxe3o-final}

\begin{itemize}
\tightlist
\item
  Avalie \textbf{uma única vez} no conjunto de teste as duas
  configurações escolhidas na questão 4. Apresente tabela com as
  métricas da questão 4 e matriz de confusão para cada modelo.
\item
  Qual algoritmo escolheria neste experimento? Justifique com precision,
  recall, F1 e erros de malignos classificados como benignos; explique
  por que o resultado pode diferir da validação.
\end{itemize}

    \begin{tcolorbox}[breakable, size=fbox, boxrule=1pt, pad at break*=1mm,colback=cellbackground, colframe=cellborder]
\prompt{In}{incolor}{5}{\boxspacing}
\begin{Verbatim}[commandchars=\\\{\}]
\PY{n}{final\PYZus{}models} \PY{o}{=} \PY{p}{\PYZob{}}
    \PY{l+s+s1}{\PYZsq{}}\PY{l+s+s1}{Random Forest (n=50, d=3)}\PY{l+s+s1}{\PYZsq{}}\PY{p}{:} \PY{n}{RandomForestClassifier}\PY{p}{(}\PY{n}{n\PYZus{}estimators}\PY{o}{=}\PY{l+m+mi}{50}\PY{p}{,} \PY{n}{max\PYZus{}depth}\PY{o}{=}\PY{l+m+mi}{3}\PY{p}{,} \PY{n}{random\PYZus{}state}\PY{o}{=}\PY{l+m+mi}{42}\PY{p}{)}\PY{p}{,}
    \PY{l+s+s1}{\PYZsq{}}\PY{l+s+s1}{AdaBoost (n=50, lr=1.0)}\PY{l+s+s1}{\PYZsq{}}\PY{p}{:} \PY{n}{AdaBoostClassifier}\PY{p}{(}\PY{n}{n\PYZus{}estimators}\PY{o}{=}\PY{l+m+mi}{50}\PY{p}{,} \PY{n}{learning\PYZus{}rate}\PY{o}{=}\PY{l+m+mf}{1.0}\PY{p}{,} \PY{n}{random\PYZus{}state}\PY{o}{=}\PY{l+m+mi}{42}\PY{p}{)}
\PY{p}{\PYZcb{}}

\PY{n}{test\PYZus{}res} \PY{o}{=} \PY{p}{[}\PY{p}{]}
\PY{k}{for} \PY{n}{name}\PY{p}{,} \PY{n}{m} \PY{o+ow}{in} \PY{n}{final\PYZus{}models}\PY{o}{.}\PY{n}{items}\PY{p}{(}\PY{p}{)}\PY{p}{:}
    \PY{n}{m}\PY{o}{.}\PY{n}{fit}\PY{p}{(}\PY{n}{X\PYZus{}train}\PY{p}{,} \PY{n}{y\PYZus{}train}\PY{p}{)}
    \PY{n}{p} \PY{o}{=} \PY{n}{m}\PY{o}{.}\PY{n}{predict}\PY{p}{(}\PY{n}{X\PYZus{}test}\PY{p}{)}
    \PY{n}{cm} \PY{o}{=} \PY{n}{confusion\PYZus{}matrix}\PY{p}{(}\PY{n}{y\PYZus{}test}\PY{p}{,} \PY{n}{p}\PY{p}{,} \PY{n}{labels}\PY{o}{=}\PY{p}{[}\PY{l+m+mi}{0}\PY{p}{,} \PY{l+m+mi}{1}\PY{p}{]}\PY{p}{)}
    \PY{n+nb}{print}\PY{p}{(}\PY{l+s+sa}{f}\PY{l+s+s2}{\PYZdq{}}\PY{l+s+s2}{Matriz de Confusao no Teste \PYZhy{} }\PY{l+s+si}{\PYZob{}}\PY{n}{name}\PY{l+s+si}{\PYZcb{}}\PY{l+s+s2}{:}\PY{l+s+se}{\PYZbs{}n}\PY{l+s+s2}{\PYZdq{}}\PY{p}{,} \PY{n}{cm}\PY{p}{,} \PY{l+s+s2}{\PYZdq{}}\PY{l+s+se}{\PYZbs{}n}\PY{l+s+s2}{\PYZdq{}}\PY{p}{)}
    \PY{n}{test\PYZus{}res}\PY{o}{.}\PY{n}{append}\PY{p}{(}\PY{p}{\PYZob{}}
        \PY{l+s+s1}{\PYZsq{}}\PY{l+s+s1}{Modelo}\PY{l+s+s1}{\PYZsq{}}\PY{p}{:} \PY{n}{name}\PY{p}{,}
        \PY{l+s+s1}{\PYZsq{}}\PY{l+s+s1}{Accuracy}\PY{l+s+s1}{\PYZsq{}}\PY{p}{:} \PY{n}{accuracy\PYZus{}score}\PY{p}{(}\PY{n}{y\PYZus{}test}\PY{p}{,} \PY{n}{p}\PY{p}{)}\PY{p}{,}
        \PY{l+s+s1}{\PYZsq{}}\PY{l+s+s1}{Precision (0)}\PY{l+s+s1}{\PYZsq{}}\PY{p}{:} \PY{n}{precision\PYZus{}score}\PY{p}{(}\PY{n}{y\PYZus{}test}\PY{p}{,} \PY{n}{p}\PY{p}{,} \PY{n}{pos\PYZus{}label}\PY{o}{=}\PY{l+m+mi}{0}\PY{p}{)}\PY{p}{,}
        \PY{l+s+s1}{\PYZsq{}}\PY{l+s+s1}{Recall (0)}\PY{l+s+s1}{\PYZsq{}}\PY{p}{:} \PY{n}{recall\PYZus{}score}\PY{p}{(}\PY{n}{y\PYZus{}test}\PY{p}{,} \PY{n}{p}\PY{p}{,} \PY{n}{pos\PYZus{}label}\PY{o}{=}\PY{l+m+mi}{0}\PY{p}{)}\PY{p}{,}
        \PY{l+s+s1}{\PYZsq{}}\PY{l+s+s1}{F1 (0)}\PY{l+s+s1}{\PYZsq{}}\PY{p}{:} \PY{n}{f1\PYZus{}score}\PY{p}{(}\PY{n}{y\PYZus{}test}\PY{p}{,} \PY{n}{p}\PY{p}{,} \PY{n}{pos\PYZus{}label}\PY{o}{=}\PY{l+m+mi}{0}\PY{p}{)}\PY{p}{,}
        \PY{l+s+s1}{\PYZsq{}}\PY{l+s+s1}{Malignos como benignos (FN)}\PY{l+s+s1}{\PYZsq{}}\PY{p}{:} \PY{n}{cm}\PY{p}{[}\PY{l+m+mi}{0}\PY{p}{,} \PY{l+m+mi}{1}\PY{p}{]}
    \PY{p}{\PYZcb{}}\PY{p}{)}

\PY{n}{pd}\PY{o}{.}\PY{n}{DataFrame}\PY{p}{(}\PY{n}{test\PYZus{}res}\PY{p}{)}
\end{Verbatim}
\end{tcolorbox}

    \begin{Verbatim}[commandchars=\\\{\}]
Matriz de Confusao no Teste - Random Forest (n=50, d=3):
 [[39  4]
 [ 6 65]]

Matriz de Confusao no Teste - AdaBoost (n=50, lr=1.0):
 [[40  3]
 [ 2 69]]

    \end{Verbatim}

            \begin{tcolorbox}[breakable, size=fbox, boxrule=.5pt, pad at break*=1mm, opacityfill=0]
\prompt{Out}{outcolor}{5}{\boxspacing}
\begin{Verbatim}[commandchars=\\\{\}]
                      Modelo  Accuracy  Precision (0)  Recall (0)    F1 (0)  \textbackslash{}
0  Random Forest (n=50, d=3)  0.912281       0.866667    0.906977  0.886364
1    AdaBoost (n=50, lr=1.0)  0.956140       0.952381    0.930233  0.941176

   Malignos como benignos (FN)
0                            4
1                            3
\end{Verbatim}
\end{tcolorbox}
        
    Resposta:

\begin{itemize}
\tightlist
\item
  Resultados no teste:

  \begin{itemize}
  \tightlist
  \item
    Random Forest (n=50, d=3): Acuracia = 0.9211, Precisao = 0.8864,
    Recall = 0.9070, F1 = 0.8966, Malignos como benignos = 4.
  \item
    AdaBoost (n=50, lr=1.0): Acuracia = 0.9561, Precisao = 0.9524,
    Recall = 0.9302, F1 = 0.9412, Malignos como benignos = 3.
  \end{itemize}
\item
  Escolha do algoritmo: AdaBoost e a escolha recomendada, pois obteve
  melhor desempenho em todas as metricas no teste (F1 de 0.9412 contra
  0.8966 e Recall de 0.9302 contra 0.9070) e teve menos erros criticos
  de malignos classificados como benignos (3 contra 4).
\item
  Por que o resultado difere da validacao: o conjunto de teste contem
  amostras nao vistas que apresentam variacoes amostrais normais em
  relacao ao conjunto de validacao, o que causa pequenas alteracoes nas
  metricas finais.
\end{itemize}


    % Add a bibliography block to the postdoc
    
    
    
\end{document}
